%% ============================================================
%% Wykaz ważniejszych oznaczeń i skrótów
%% ============================================================

%% ------------------ Nagłówek + wpis do spisu treści -------------------
\chapter*{Wykaz ważniejszych oznaczeń i skrótów}
\addcontentsline{toc}{chapter}{Wykaz ważniejszych oznaczeń i skrótów}

%% --------- Podstawowe wielkości i operatory ----------
%% Tabela 4 kolumny: symbol | jednostka | łącznik „--” | opis
%% Wyrównanie: l r c l (lewo, prawo, środek, lewo)
\section*{Podstawowe wielkości fizyczne oraz wyrażenia matematyczne}
\noindent
\begin{tabular}{l r c l}
     $j$  &$[\cdot]$  &--&jednostka urojona\\
	 $\nabla$  &$[\cdot]$  &--&operator nabla\\
	 $\vec{i}_{x},\vec{i}_{z},\vec{i}_{y} $ &$[\cdot]$  &--& wersory układu współrzędnych kartezjańskich    \\
	 $\vec{F}$ &$[\text{N}]$  &--& siła \\
	 $t $ &$[\text{s}]$  &--& czas \\
	 $m$  &$[\text{kg}]$  &--& masa \\
	 $R$  &$[\Omega]$  &--&rezystancja\\
	 $\mathbf{S}$  &$[\cdot]$  &--&macierz rozproszenia\\
\end{tabular}

%% -------------- Stałe i wartości znane ---------------
%% Tabela 4 kolumny: symbol/wartość | jednostka | łącznik | opis
\section*{Stałe fizyczne i matematyczne}
\noindent
	 \begin{tabular}{l r c l}
	$\pi \approx 3.14159 $&$[\cdot]  $ &--& liczba $\pi$\\ 
	$g \approx 9.80665 $&$[\frac{\text{m}}{\text{s}^2}]  $ &--& przyspieszenie ziemskie\\
\end{tabular}

%% ---------------- Terminologia użyta w pracy -----------------
%% Tabela 3 kolumny: komenda | łącznik „--” | objaśnienie
%% Użyte \verb|...| zabezpiecza zapis komend dosłownie.
%% \multirow rozszerza opis na kilka wierszy
\section*{Terminologia użyta w pracy projektowej}
\noindent
\begin{tabular}{l c l} 
\verb|\section{}| &--& tworzy tytuł rozdziału \\
\verb|\subsection{}|&--&  tworzy tytuł podrozdziału \\
\verb|\cite{}| &--& odwołanie do pozycji z bibliografii \\
\verb|\ref{}|&--&  odwołanie do rysunku, tabeli lub wzoru \\
\verb|\eqref{}| &--& odwołanie do równania (z nawiasem)  \\
\verb|\begin{figure}| & \multirow{3}{*}{--} & \multirow{3}{*}{wstawienie obiektu graficznego} \\
... & & \\ 
\verb|end{figure}| & & \\ 
\verb|\begin{table}| & \multirow{3}{*}{--} & \multirow{3}{*}{wstawienie tabeli} \\
... & & \\ 
\verb|end{table}| & & \\ 
\verb|\begin{equation}| & \multirow{3}{*}{--} & \multirow{3}{*}{wstawienie wzoru z numeracją} \\
... & & \\ 
\verb|end{equation}| & & \\ 
\verb|\textbf{}| &--& pogrubienie tekstu  \\
\verb|\emph{}| &--& wyróżnienie kursywą  \\
\verb|\footnote{}| &--& wstawienie przypisu dolnego \\
\verb|\includegraphics{}| &--& wstawienie obrazu \\
\verb|\label{}| &--& etykieta obiektu \\
\verb|\item| &--& element listy \\
\end{tabular}
